\documentclass[a4paper, 12pt]{article}

% URL line-breaking (must precede any package that loads url)
\PassOptionsToPackage{hyphens}{url}

% Encoding & font
\usepackage[utf8]{inputenc}
\usepackage[T1]{fontenc}

% Page layout
\usepackage[a4paper, margin=2.5cm]{geometry}


% Mathematics
\usepackage{amsmath}
\usepackage{amssymb}
\usepackage{amsthm}

\newtheorem{theorem}{Theorem}[section]
\newtheorem{lemma}[theorem]{Lemma}
\newtheorem{definition}[theorem]{Definition}
\newtheorem{corollary}[theorem]{Corollary}

% Graphics & figures
\usepackage{graphicx}
\usepackage{float}              % [H] placement
\usepackage{caption}
\usepackage{subcaption}

% Code listings
\usepackage{listings}
\usepackage{xcolor}

% References & hyperlinks
\usepackage[
  backend  = biber,
  style    = numeric,
  sorting  = none,
]{biblatex}
\addbibresource{references.bib}

\usepackage{hyperref}
\hypersetup{
  colorlinks = true,
  linkcolor  = black,
  citecolor  = black,
  urlcolor   = blue,
}

% syntax highlighting for code examples
\usepackage{minted}

% Miscellaneous
\usepackage{todonotes}          % \todo{...} margin notes during drafting


\title{Analysing the Treewidth of Rust Programs}
\author{
  Battini, Liam
  \and
  Haans, Geert
  \and
  Singh, Prabhjot
  \and
  Steffens, Tammo
}

\date{\today}

\begin{document}

\maketitle

\begin{abstract}
	Treewidth is a graph-theoretic measure of structural complexity with direct applications in compiler optimisation.
	Thorup~(1998) showed that the control-flow graphs~(CFGs) of goto-free C programs have treewidth at most~6, enabling efficient register allocation.
	Gustedt, M{\ae}hle, and Telle~(2002) extended this result to label-free Java and confirmed it empirically on the Java standard library, finding an average treewidth of~2.7 and a maximum of~5.
	In this paper, we perform an analogous empirical study for Rust.
	Rather than analysing the source language, we operate on Rust's Mid-level Intermediate Representation~(MIR), the compiler's own explicit CFG-based representation at which flow-sensitive analyses such as borrow checking are performed.
	We extract CFGs for all functions in the Rust standard library using a custom compiler driver and compute their treewidth with \texttt{FlowCutter}.
	\todo[inline]{Add results sentence: mean, median, max treewidth; safe vs.\ unsafe split.}
	We report results for all functions combined and separately for safe and unsafe Rust.
\end{abstract}

\section{Introduction}
\label{sec:introduction}

Control-flow graphs (CFGs) are a fundamental abstraction used by compilers to represent the execution structure of programs. Many optimisation and analysis tasks in compilers operate directly on these graphs~\cite{allen_1970}. A key structural property of graphs is treewidth, a parameter that measures how close a graph is to a tree. The treewidth is obtained by the tree decomposition, which is a mapping of a graph into a tree, where the vertices are a set of nodes from the original graph. The treewidth is the size of the largest vertex (or bag) minus one~\cite{bodlaender_1993}. Many NP-hard problems can algorithmically be solved in linear time when the graph has a smaller treewidth~\cite{thorup_1998}. This observation makes the treewidth of CFGs an important question for compiler design.

Register allocation is an example of an optimisation problem in compilers, which can be represented as a graph problem, where the complexity depends on the structure of the underlying program representation~\cite{thorup_1998}. If the CFG of the program has a bounded treewidth, dynamic-programming can be used to compute the optimal allocations efficiently. Understanding the structural complexity of real program CFGs therefore has direct implications for the feasibility of exact optimisation techniques in compilers~\cite{bodlaender_1988}.

The connection between structured programs (goto-free) and bounded treewidth was shown in~\cite{thorup_1998}, it was shown that the CFGs of structured programs have a bounded treewidth. In particular, he proved that CFGs of goto-free C programs have a treewidth of at most six. This finding implies that the structural complexity of control flow is limited by structured programming structures. Building on this observation,~\cite{gustedt_2002} investigated whether the same behaviour happens in practice for Java programs. It was shown that label-free Java has a bound of at most six and confirmed that real Java methods exhibit small treewidth values, with an average of 2.7 and a maximum of five in the Java standard library.

Based on these result, we raise a similar question for modern programming languages: \textit{what is the treewidth of Rust programs?} Rather than analysing source-level control flow, we operate on the Mid-level Intermediate Representation (MIR), the compiler's internal CFG representation on which our measurements are based.

By intercepting the Rust compiler pipeline after the MIR optimisation, we can extract the CFG of each function. We then compute the treewidth of each extracted CFG using FlowCutter. This allows us to measure the structural complexity of a Rust program, at the same level as where the compiler performs its analyses.

Our contributions are as following:
\begin{enumerate}
	\item \textbf{Theoretical:} We study Rust within the existing work on treewidth of programming languages, by analysing Rust's control-flow constructs to prior results for structured languages such as Java.
	\item \textbf{Tool/Emperical:} We develop a toolchain for extracting CFGs from the Rust MIR and computing the treewidth.
\end{enumerate}

This paper is organised as follows: In section 2, we introduce the necessary preliminaries and background concepts, namely Control-Flow graphs (CFGs) and treewidth, where we formally defined the tree decomposition. Section 3 reviews some of the related work for our research, for example the treewidth of Java programs. Section 4 provides background on Rust and its Mid-level Intermediate Representation (MIR). Followed by section 5, which describes the methodology for extracting CFGs and computing the treewidth. Section 6 presents the results of our analysis. Finally, in section 7 we summarize the results and outline directions for future work.

% - Motivation: why treewidth of programs matters
% - Connection to register allocation and compiler optimisation
% - Thorup's result for goto-free C; Gustedt et al. for Java
% - Research question: what about Rust?
% - Contributions: (1) theoretical, (2) tool/empirical
% - Outline of the paper

\section{Background}
\label{sec:background}

\subsection{Control-Flow Graphs}
\label{subsec:cfg}

Thorup~\cite{thorup_1998} states control-flow graphs (CFGs) are graphs where nodes represent basic blocks and edges represent possible control flow paths between them. According to Allen~\cite{control_flow_analysis}, basic blocks are straight-line sequence of instructions that have single entry and exit points, where inside a block no jumps are allowed and disallows jumps from outside the block into the middle of the block. The edges in a CFG sets the constraint on how execution can flow between these blocks. 

The topology of CFGs is determined by program's control structure. Sequential statements produce linear chains of basic blocks, while additional control constructs such as conditionals, loops and function exits makes the flow between nodes more complicated. For example, loops creates edges that point back to earlier blocks in the graph. 

% - Definition: nodes = basic blocks, edges = control transfers
% - Relationship between program structure and CFG topology
% - CFG treewidth as a proxy for program complexity

\subsection{Treewidth}
\label{subsec:treewidth}
Treewidth is a measurement on how "tree-like" a graph is and it can be measured on any graph. This is useful because problems that are more "tree-like" can be solved faster, meaning a lower treewidth means it is quicker to solve the problem. 

\subsubsection{Tree decomposition}
In order to measure the treewidth of a graph, a tree decomposition must first be created from the given graph. This can be done by representing the original graph \texttt{G = (V,E)} in a tree \texttt{T} with overlapping sets of nodes, known as bags. The constraints for a tree decomposition, is that every node in the graph must be covered in at least one bag which gurantees every node is covered in th

% - Definition of treewidth (tree decomposition)
% - Relevance: bounded treewidth enables efficient dynamic programming
% - Known bounds for structured programs (Thorup, Gustedt et al.)
\section{Related Work}
\label{sec:related-work}

\subsection{Treewidth of Structured Programs}
\label{subsec:rw-treewidth-structured}

% - Thorup (1998): treewidth \leq 6 for goto-free C; linear-time register allocation
% - Bodlaender, Gustedt, Telle (1998): linear-time register allocation for fixed number of registers

\subsection{Treewidth of Java Programs}
\label{subsec:rw-java}

Gustedt, Mæhle, and Telle~\cite{gustedt_treewidth_2002} investigate whether Thorup's result for \texttt{goto}-free C~\cite{thorup_1998} carries over to Java.
Thorup showed that \texttt{goto}-free C programs have CFG treewidth at most~6, using the notion of \emph{flow-affecting constructs} (FACs): language constructs that add non-local control-flow edges and raise treewidth by at most~1 each.

Java has no \texttt{goto}, but its labeled \texttt{break} and \texttt{continue} statements are expressively equivalent to \texttt{goto}.
Consequently, full Java programs can have arbitrarily high treewidth.
Restricting to a label-free subset eliminates this issue: label-free Java has exactly four FACs (\texttt{break}, \texttt{continue}, \texttt{return}, and \texttt{short-circuit evaluation}), giving a treewidth bound of at most~$2 + 4 = 6$.

To complement this theoretical result, the authors implement a tool that computes tree decompositions of Java CFGs and apply it to 9387 methods from the Java standard library.
The average treewidth is~2.7 and the maximum observed is~5, with over 95\% of methods having treewidth at most~3, which is well below the theoretical bound.

This paper is directly relevant to our work: it establishes both a theoretical treewidth bound and empirical baselines for a real-world language, serving as the primary comparison point for our analogous analysis of Rust programs via MIR.

\subsection{Static Analysis of Rust}
\label{subsec:rw-rust-analysis}

% - Prusti: formal verification via MIR
% - Budde (2024): data-flow vulnerability analysis; CFG extraction tooling we may reuse
% - Other MIR-based tools

\section{Rust and MIR}
\label{sec:rust-and-mir}

This section briefly situates Rust relative to Java~\cite{gustedt_2002}, describes the features of Rust relevant to treewidth, and introduces MIR, the compiler's internal CFG representation on which our measurements are based.
Rust is a systems programming language designed to combine low-level control over memory with strong compile-time safety guarantees.
Unlike Java, Rust does not rely on a garbage collector; instead, memory safety is enforced statically at compile time by the \emph{borrow checker}~\cite[Ch.~4.2]{rust_book}.

\subsection{Ownership and Borrowing}
\label{subsec:ownership-borrowing}

Every value in Rust has a unique \emph{owner}~\cite[Ch.~4.1]{rust_book}.
Ownership can be transferred by assignment or by passing a value to a function.
After such a transfer, the original variable may no longer be used, and when the owner goes out of scope, the value is automatically deallocated.
This prevents multiple independent variables from being responsible for the same piece of memory.

Programs sometimes need temporary access to a value without transferring ownership.
Rust supports this through \emph{borrowing}~\cite[Ch.~4.2]{rust_book}: a borrow is a reference to a value owned elsewhere.
There are two kinds: \emph{shared borrows} (\texttt{\&T}), which allow read-only access, and \emph{mutable borrows} (\texttt{\&mut T}), which allow modification.
Rust enforces strict aliasing rules: any number of shared borrows may coexist, or exactly one mutable borrow may exist, but not both at the same time.
A borrow is only valid for some region of the program, called its \emph{lifetime}~\cite[Ch.~10.3]{rust_book}.
References may not outlive the values they refer to, preventing dangling pointers statically.
Together, these aliasing and lifetime constraints are part of Rust's static semantics, enforced at compile time by the borrow checker.
These checks are \emph{flow-sensitive}: whether a borrow is active depends on the control-flow path taken through the program.
The borrow checker therefore operates directly on the CFG, and ownership constraints can influence its structure: values that go out of scope on some paths but not others require explicit drop edges in the compiled control-flow graph.

\subsection{Rust Control Flow}
\label{subsec:rust-control-flow}

The structure of Rust's control flow is the other main factor shaping CFG complexity.
Rust uses predominantly \emph{structured} control flow: execution is expressed through conditional expressions, pattern matching, and loops rather than arbitrary jumps~\cite[Ch. 8.2.13]{rust_reference}.
This is one of the main reasons CFGs of real programs tend to remain relatively simple, as prior work on goto-free language fragments has shown~\cite{thorup_1998}.

Certain language features can make the compiled control-flow graph more complex than the source structure suggests.
Labeled \texttt{break} and \texttt{continue} can jump to an enclosing loop, creating edges in the CFG that bypass intermediate blocks.
Gustedt et al.\ showed that such constructs are similarly harmful to \texttt{goto} and can lead to arbitrarily high treewidth~\cite{gustedt_2002}.
Similarly, async functions compile to state machines whose branches and cycles are not visible at the source level.
This motivates measuring treewidth empirically on the compiler's internal representation (see Section~\ref{subsec:mir}) rather than reasoning from source-level nesting alone.

\subsection{Mid-level Intermediate Representation}
\label{subsec:mir}

The Rust compiler lowers functions into the \emph{Mid-level Intermediate Representation} (MIR) after parsing and earlier compilation stages~\cite[Ch.~44]{rustc_dev_guide}.
MIR is a control-flow graph: it consists of basic blocks containing simple typed statements, together with explicit control-flow transfers between those blocks.
A basic block is a sequence of statements followed by a \emph{terminator}.
Terminators represent control transfers with potentially multiple successors, covering branching, function calls, drops, and returns.
Once lowered to MIR, execution order is no longer implicit in the source structure but is represented explicitly as a graph.

MIR provides a direct basis for studying the structure of real Rust programs.
It offers a CFG representation suitable for treewidth analysis, and coincides with the level at which the compiler performs its own flow-sensitive analyses~\cite[Ch.~44]{rustc_dev_guide}.
We therefore extract and analyse CFGs at the MIR level, where the compiler's own structural reasoning takes place.
The next section covers this in detail.

\section{Approach}
\label{sec:approach}

\subsection{Rust Subset for Theoretical Analysis}
\label{subsec:rust-subset}

% - Define a structured subset of Rust analogous to goto-free C / label-free Java
% - Discuss constructs excluded and why
% - Relationship to the borrow checker

\subsection{CFG Extraction from MIR}
\label{subsec:cfg-extraction}

% - Overview of the extraction pipeline
% - Tooling: reuse/adaptation of Budde (2024) data-flow tool
% - Mapping MIR basic blocks to CFG nodes/edges

\subsection{Treewidth Computation}
\label{subsec:treewidth-computation}

% - Algorithm used for computing treewidth (e.g. exact vs. heuristic)
% - Implementation details

\section{Evaluation}
\label{sec:evaluation}

\subsection{Experimental Setup}
\label{subsec:setup}

Using the pipeline from Section~\ref{sec:methodology}, we analysed the Rust standard library for the host target \texttt{aarch64-apple-darwin}. The corpus consists of the \texttt{core}, \texttt{alloc}, and \texttt{std} crates compiled with monomorphisation enabled, so each concrete instantiation of a generic function contributes its own MIR body and CFG. In total, the corpus contains 27{,}997 functions: 22{,}273 safe functions and 5{,}724 functions declared as \texttt{unsafe fn}.

For each function, we computed the treewidth of its MIR CFG. We report the mean, median, and maximum treewidth, together with the share of functions with treewidth at most~3 and at most~6. The threshold of~6 is included to facilitate comparison with the bound established for goto-free C and for label-free Java~\cite{thorup_1998,gustedt_2002}.

\subsection{Results}
\label{subsec:results}

Treewidth is very small across the corpus. Over all 27{,}997 functions, the mean treewidth is 1.04, the median is 1, and the maximum observed treewidth is 6. Moreover, 99.65\% of all functions have treewidth at most~3, and all functions have treewidth at most~6. Thus, MIR CFGs in the Rust standard library are overwhelmingly close to trees in practice.

\begin{table}[h]
	\centering
	\begin{tabular}{lrrrrr}
		\hline
		Group            & Count    & Mean & Median & Max & $\mathrm{tw} \leq 3$ \\
		\hline
		All functions    & 27{,}997 & 1.04 & 1      & 6   & 99.65\%              \\
		Safe functions   & 22{,}273 & 1.06 & 1      & 6   & 99.59\%              \\
		Unsafe functions & 5{,}724  & 0.97 & 1      & 4   & 99.90\%              \\
		\hline
	\end{tabular}
	\caption{Treewidth summary for the Rust standard library MIR corpus.}
	\label{tab:evaluation-summary}
\end{table}

\noindent Table~\ref{tab:evaluation-distribution} shows that the distribution is strongly concentrated at very small values: 97.90\% of all functions have treewidth at most~2, and only 97 functions (0.35\%) exceed treewidth~3. The largest values occur in a small number of formatting and systems-oriented routines with unusually branch-heavy control flow.
\begin{table}[h]
	\centering
	\begin{tabular}{lrr}
		\hline
		Treewidth & Count  & Share \\
		\hline
		0         & 3{,}853 & 13.76\% \\
		1         & 19{,}773 & 70.63\% \\
		2         & 3{,}783 & 13.51\% \\
		3         & 491    & 1.75\% \\
		4         & 90     & 0.32\% \\
		5         & 5      & 0.02\% \\
		6         & 2      & 0.01\% \\
		\hline
	\end{tabular}
    \caption{Distribution of treewidth values across all functions in the Rust standard library
   MIR corpus.}
    \label{tab:evaluation-distribution}
\end{table}

The safe/unsafe split does not indicate higher structural complexity for unsafe Rust in this corpus. Unsafe functions have a slightly lower mean treewidth than safe functions (0.97 versus 1.06), and their maximum observed treewidth is only~4.

\subsection{Discussion}
\label{subsec:discussion}

Overall the results support the hypothesis that practical Rust CFGs have small treewidth. Even without a Rust-specific structural proof analogous to Thorup's result~\cite{thorup_1998}, the empirical picture is clear: almost all MIR CFGs lie well below the bound of~6 associated with structured C and Java fragments, and the maximum observed value is exactly~6.

For further context, the results invite comparison with the Java study of Gustedt et al., which should be interpreted with care: their analysis is source-level and structurally constrained, while ours is based on MIR basic-block CFGs~\cite{gustedt_2002}.
In particular, treewidth~0 and~1 are common in our setting but cannot occur in their construction due to Thorup's decomposition algorithm.
With this in mind, they report an average treewidth of about~2.7 and a maximum of~5 for the Java standard library, whereas we obtain a mean of~1.04 and a maximum of~6 for Rust MIR.

While these results are promising, several limitations should be considered.
First, the study is limited to the Rust standard library and may not represent application code or third-party crates.
Second, the analysis is performed on optimised MIR for a single compiler version and host target, so the exact CFGs may vary across releases and platforms.
Third, we analyse undirected CFGs after removing unwind-only cleanup edges, which matches the graph-theoretic notion of treewidth but abstracts away some aspects of Rust's exceptional control flow.

Despite these limitations, the results provide a consistent empirical baseline for the structural complexity of Rust MIR CFGs.
The consistently small treewidth values suggest that compiler optimisations designed for graphs of bounded treewidth are well-suited to real Rust MIR in practice.

\section{Conclusion}
\label{sec:conclusion}

% - Summary of contributions
% - Answer to the research question
% - Limitations and future work


\printbibliography

\newpage
\section*{AI Statement}

% Describe in detail how AI tools (e.g.\ ChatGPT, GitHub Copilot, or similar)
% were used during the research and writing of this paper.
% If no AI tools were used, a brief statement to that effect is sufficient.

\end{document}
